% Generated draft: requires mathematical and source review.
\documentclass[11pt]{article}
\usepackage[margin=25mm]{geometry}
\usepackage{fontspec}
\setmainfont{Latin Modern Roman}
\usepackage{amsmath,amssymb,mathtools,unicode-math}
\setmathfont{Latin Modern Math}
\usepackage{enumitem,needspace,xurl,hyperref}
\setlength{\parindent}{0pt}
\setlength{\parskip}{5pt}
\setlength{\emergencystretch}{3em}
\newcommand{\N}{\mathbb{N}}
\newcommand{\Z}{\mathbb{Z}}
\newcommand{\Q}{\mathbb{Q}}
\newcommand{\R}{\mathbb{R}}
\newcommand{\C}{\mathbb{C}}
\newcommand{\F}{\mathbb{F}}
\newcommand{\A}{\mathbb{A}}
\newcommand{\PP}{\mathbb{P}}
\newcommand{\E}{\mathbb{E}}
\newcommand{\eps}{\varepsilon}
\newcommand{\dd}{\,\mathrm{d}}
\begin{document}
\section*{UnsolvedMath ID: 30003963}
\textbf{Source problem code:} catalogue entry
\subsection{Definitions and mathematical statement}
The source specifies a degree \(d\), but does not specify the genus, number of markings, base field or base scheme, or category of moduli spaces. The following notation describes the conventional interpretation conditionally on those choices; it does not select a particular compactification.

Fix integers \(g,n\geq 0\) with \(2g-2+n>0\), and \(d\in\Z\). After choosing a base, let \(\mathcal M_{g,n}\) denote the moduli stack of smooth, proper, geometrically connected genus-\(g\) curves with \(n\) ordered, distinct marked points. Let \(\overline{\mathcal M}_{g,n}\) denote the moduli stack of stable marked curves: the curves may have nodes, the markings are distinct smooth points, and the marked curves have finite automorphism groups. Write \(j:\mathcal M_{g,n}\hookrightarrow\overline{\mathcal M}_{g,n}\) for the smooth locus. Write \(\mathcal J^d_{g,n}\to\mathcal M_{g,n}\) schematically for the degree-\(d\) universal Jacobian moduli problem, whose fiber over a smooth marked curve \(C\) is \(\operatorname{Pic}^d(C)\), parametrizing isomorphism classes of degree-\(d\) line bundles. A precise construction must specify whether this notation denotes a relative Picard space, a rigidified stack, or the stack retaining scalar automorphisms of line bundles.

For a geometric stable curve \(C\), a possible replacement for a line bundle is a coherent sheaf \(F\) that is torsion-free and has rank one at the generic point of every irreducible component. Its total degree is
\[
\deg_C(F)=\chi(C,F)-\chi(C,\mathcal O_C).
\]
If \(C_1,\ldots,C_r\) are the irreducible components of \(C\), the multidegree of a line bundle \(L\) is
\[
\underline{\deg}(L)=\bigl(\deg_{C_1}(L|_{C_1}),\ldots,\deg_{C_r}(L|_{C_r})\bigr),\qquad
\sum_{i=1}^{r}\deg_{C_i}(L|_{C_i})=d.
\]
For a torsion-free sheaf, component degrees can instead be defined using the restrictions modulo torsion; their relation to total degree must account for failure of local freeness at nodes. The source does not prescribe a sheaf model or a multidegree convention for such boundary objects.

\textbf{Problem.} Determine how to construct an extension
\[
\overline\pi:\overline{\mathcal J}^{\,d}_{g,n}\longrightarrow\overline{\mathcal M}_{g,n},
\qquad
\overline{\mathcal J}^{\,d}_{g,n}\times_{\overline{\mathcal M}_{g,n}}\mathcal M_{g,n}
\simeq\mathcal J^d_{g,n},
\]
whose fibers over stable marked curves are well-behaved compactified Jacobians, and determine which degenerations of line bundles and which multidegrees belong in these fibers. Specify the auxiliary choices and conditions needed to make this construction precise, including any stability condition or polarization used to select boundary objects.

The phrase ``well-behaved'' is not defined in the supplied record. In particular, the record does not specify the required properness, separatedness, flatness, representability, universal-family, or base-change properties, nor the precise locus that must be dense in a compactified fiber. Consequently, this is a formulation of the supplied construction question with explicit unresolved parameters, rather than a uniquely specified existence assertion.
\subsection{Short English statement}
Extend the degree-\(d\) universal Jacobian from smooth marked curves to stable marked curves, and determine the boundary objects and multidegrees to include. The source leaves the base, moduli category, stability choices, and meaning of ``well-behaved'' unspecified, so these remain part of the formulation requiring review.
\subsection{Sources}
\begin{itemize}
\item Local UnsolvedMath registry, supplied complete record: permanent ID \(30003963\), title ``Compactified Universal Jacobians over Stable Curves,'' exact source code \texttt{catalogue entry}, collection ``Oberwolfach Reports Open Problems.'' No individual proposer is supplied; the proposed year is recorded as 2018.
\item Supplied source citation: \emph{Flat Surfaces and Algebraic Curves} (2019), Oberwolfach Reports, DOI \texttt{10.4171/owr/2018/42}, pp.\ 2583--2650. \url{https://doi.org/10.4171/owr/2018/42}. The background attributes the question to workshop open-problem discussions documented in Oberwolfach Reports. No precise page or passage for this question is supplied.
\item Supplied literature source: ``Oberwolfach Report 42, compactified Jacobians,'' \url{https://doi.org/10.4171/owr/2018/42}. Literature assessment checked at \texttt{2026-08-21}: this is supplied dated curation, not independently verified. It labels the problem partially solved and reports constructions of compactified Jacobians and proper moduli of stable torsion-free sheaves for broad stability data, together with wall-crossing descriptions, while reporting no single complete degeneration theory covering all curves and choices. The background expressly requests independent verification of time-sensitive status and cited claims.
\item The registry labels the statement \texttt{corrected\_verified} and reports manual semantic, typographic, and LaTeX review. These are supplied registry assertions, not independent verification of the source passage or mathematical correctness of this draft.
\end{itemize}
\end{document}
